\section{Towards a Quantum Pseudocode}\label{sec:pseudocode}
  \subsection{Quantum Entities and Hidden Entanglement}
    In \autoref{sec:primer}, we defined notation for computational basis states
      $\ket{x}$, general quantum states $\ket{\psi}$,
      unitary transformations $\unitary$, and the
      the result of applying a unitary transformation $\unitary$ on a
      quantum state $\ket{\psi}$, written 
      $\ket{\psi} \apply{\unitary} \ket{\varphi}$; see \autoref{tab:notation}
      for an overview.

    And still this isn't enough. This is because we are going to be talking
      about communicating parties. Consider the situation that Alice has two
      quantum bits initalized to the zero state, so that her state is
      $\ket{\psi} = \ket{0}_0 \ket{0}_1$. Now she wants to send her second
      qubit register, $\ket{0}_1$, over to Bob, who is going to apply some
      single-qubit unitary to it, and then send it back. We might write this as
      $\ket{0}_1 \apply{\unitary_1} \ket{\xi}_1$, where $\ket{\xi}_1$ is some
      unknown, single-qubit state, and say that Bob returns $\ket{\xi}_1$ to
      Alice at the end of the protocol.

    But what if Alice had prepared her two qubits in the state
      $\ket{\varphi} = \frac{\ket{00} + \ket{11}}{\sqrt{2}}$ instead?
      Then, the two-qubit state is \emph{non-separable}, meaning that it can no
      longer be cleanly factored into a tensor product of two single-qubit
      states. In the world of quantum
      physics, this means that the two qubits are \emph{entangled}---a fancy
      word describing a phenomenon with many deep consequences for physics and
      philosophy alike---but mathematically speaking it is nothing but the quantum
      generalization of correlated probabilities.\footnote{
        In other words, the term ``separable'' here carries exactly the same
        meaning as the term ``product'' does when talking about classical
        probability distributions.
        More precisely, if an arbitrary quantum state is \emph{separable} along
        some dividing line, such that $\ket{\psi} = \ket{\varphi} \otimes
        \ket{\xi}$, then the probability distribution over the possible
        outcomes of observing $\ket{\psi}$ will
        form a product distribution, factorizable along the
        same dividing line. This remains true regardless of
        which basis you use to observe the state, implying that
        entanglement is an observation-independent (and therefore
        ``objective'') property of the state.
      } Alice can still send her second qubit to Bob, and Bob can still apply
      the operation $\unitary$ to it; our problem is just that with the tools
      presented so far, we have no way\footnote{
        For the readers ahead of the curve: Yes, one may reference one half of an
        entangled state using the generalized notion of a 
        \emph{mixed state}, by going to the density matrix formalization of
        quantum mechanics, which generalizes quantum states by combining them
        with classical probabilities. (Informally speaking, a density matrix is
        nothing but a probability distribution over quantum states.) The reason
        we avoid this solution is that mixed states are not
        \emph{unique}, in the sense that there are many (larger) pure
        states that might have given rise to a particular density matrix, 
        and so there would be many global states that could have led to Bob
        receiving a particular mixed state. Therefore, density matrices lack
        the expressive power to specify how Bob's operations on the qubits
        available to him may, through entanglement, affect the state of Alice's qubits.
      } of writing the operation that Bob is
      performing on $\ket{0}_1$ without referencing the entire
      state, like so:
      \[
        \frac{\ket{00}_{0,1} + \ket{11}_{0,1}}{\sqrt{2}} \longrightarrow \left( 
        \idmatrix_0 \otimes \unitary_1 \right) 
        \left( \frac{\ket{00}_{0,1} + \ket{11}_{0,1}}{\sqrt{2}} \right)\, .
      \] 

    \noindent
    What we want, is to be able to write the operations of Alice and Bob without
      having to reference qubits that they don't have access to, regardless of whether
      or not their quantum states are entangled. In other words, what we need is a
      better \emph{quantum pseudocode}.

  \subsection{Qubits and Registers}
    Let $\qubit$ be a physical qubit, in the sense of being some kind of
      particle or quantum system existing in
      the physical world. Now, guided by our intuition, we want to say that Alice,
      instead of sending a quantum \emph{state} to Bob, actually sends him the
      \emph{physical qubit} $\qubit$.\footnote{
        In real-world quantum technologies, qubits are
        typically stored in one kind of system (e.g.\ neutral atoms) and sent using another
        kind (usually photons), but due to the non-cloneability of quantum states,
        together with the universal ability to transfer any quantum state from
        one qubit to another via the $\SWAPgate$ gate, this simplified
        model is without loss of generality.
      } Anyone who receives a qubit, should also gain the
      ability to apply any transformation or measurement on it, while those who sent it 
      should lose the ability to do so, unless and until the qubit is returned to them.
      Crucially, we should be able to describe these actions \emph{locally},
      i.e.\ without reference to any qubits \emph{not} in the possession of the
      receiving party.

    So, with that goal in mind, let us introduce a little more formalism.

    \begin{definition}[Qubit]\label{def:qubit}
      A qubit $\qubit$ is a two-dimensional physical system, i.e.\ a physical
      system associated with an \emph{observable} $\hat{\qubit}$, taking values
      in $\bin$. Each qubit is additionally associated with 
      a unique \emph{register address} $i \in \NN$.\footnote{
        In the formalism of quantum Turing machines given by
        Deutsch~\cite{Deutsch85}, the addresses correspond
        to values of the \emph{address observable} $\hat{x}$
        (except that we will limit ourselves to finite quantum
        memory by letting the set of possible addresses be $[n]$
        rather than $\ZZ$).
      }     
    \end{definition}

    \noindent
    Let us refer to the qubit sitting at address $i$ as $\qubit_i$.
      If qubit $\qubit_i$ is in not entangled with any
      other qubit (i.e.\ is in a pure state), then its state may be denoted
      $\ket{\qubit_i}$ (equivalent to $\ket{\qubit_i}_i$). If
      $\ket{\qubit_i} = \ket{\psi}$ for some particular single-qubit
      state $\ket{\psi}$, then this may be denoted $\ket{\psi}_i$. Likewise,
      the state of qubit $i$ at time $t$ may be denoted $\ket{\qubit_i(t)}$.
      For the sake of illustration, let's say that at
      time $t$ the unitary $\unitary$ was applied to $\qubit_i$; this would then be written
      $\ket{\qubit_i(t)} \apply{\unitary_i} \ket{\qubit_i(t+1)}$. Likewise, if $\qubit$ was
      initialized to $\ket0$, then $\ket{\qubit(0)} = \ket0$.

    Next, let us define our ``universe'', that is, the full, undivided quantum memory 
      which will act as a stage for all our plays.

    \begin{definition}[Quantum memory]\label{def:memory}
      An $n$-qubit quantum memory is an $n$-element list of qubits $\qcomp =
      (\qubit_1, \cdots, \qubit_n)$ comprising a quantum
      system that is \emph{closed} in the following sense: 
      There are no operations that may be performed on the qubits $\qubit_i$
      such that $\qcomp$ can no longer be said to be in a pure state.
    \end{definition}

    \noindent
    In particular, $\qcomp$ can not get entangled with any
      system outside of itself---though it may be \emph{observed}, after which
      it will immediately collapse to the observed basis state.\footnote{
        Of course, from the Everettian point of view, observing a state
        \emph{is} to entangle yourself with it.
      }         

    The requirement that $\qcomp$ is always in a pure state also
      implies that state preparation must be
      \emph{deterministic} (or at least posterior to any randomness
      involved in the procedure): Flipping a (classical) coin and saying that you
      prepare a qubit in the state $\ket{1}$ on heads and $\ket{0}$ on
      tails would, when conditioning on the result of the coin flip, put $\qcomp$
      in a mixed state, and this must therefore be disallowed. Thankfully, due to
      Turing-universality, we can make this restriction without loss of generality.

    Thus, our $N$-dimensional ``global'' quantum state $\ket{\qcomp}$ will always
      be well-defined, giving us a ``smallest'' element $\qubit$ and a
      ``largest'' element $\qcomp$. That means it is time to define our in-betweens.

    \begin{definition}[Quantum register]\label{def:register}
      An $m$-qubit quantum register $\register$ is a subsystem of the $n$-qubit quantum
      memory $\qcomp$, that is, an ordered sub-list of $\qcomp$, where by
      \emph{ordered} we mean that if
      $\register = (\qubit_{i_0}, \qubit_{i_1}, \dots, \qubit_{i_{m-1}})$
      then $i_0 < i_1 < \dots < i_{m-1}$.
    \end{definition}

    \noindent
    If the qubits in $\register$ are not entangled with any other subsystems
      of $\qcomp$, then the system $\register$ is in a pure state, and we may
      write $\ket{\register}$. In a slight overload of notation, let
      $[\register]$ denote the set $\{i_j\}_{j\in [m]}$ of qubit addresses
      associated with the qubits in $\register$; then $\ket\register$ is given
      by
    \[ 
      \ket\register 
      \coloneqq \bigotimes_{\forall i \in [m]} \ket{\register[A][i]}
      \coloneqq \bigotimes_{\forall i \in [\register]} \ket{\qcomp[i]}
      = \ket{\qubit_{i_0}, \ldots, \qubit_{i_{m-1}}}\, .
    \]

    \noindent
    If $\ket{\register} = \ket{\psi}$ for some specific state $\ket{\psi}$,
      this may be denoted $\ket{\psi}_{[\register]}$,
      which, having satisfied ourselves that it is
      well-defined, we will immediately shorthand to $\ket{\psi}_{\register}$.

  \subsection{Operating on Registers}
    But of course, the starting point of this entire exercise was that we
      \emph{cannot} write
      the state of a register in terms of its own constituents in the general,
      entangled case, and so we still need an alternative way to denote 
      operating on it.
    
    \begin{definition}[Operating on registers]\label{def:apply}
      Let $\qcomp$ be an $n$-qubit quantum memory in quantum state
      $\ket{\qcomp}$, let the register $\register$ be an $m$-qubit
      subsystem of $\qcomp$, and let $\unitary$ be an
      $M \times M$ unitary matrix. 

      Next, define $\unitary_{\register}$ to be the $N \times N$ unitary matrix that
      has the effect of applying $\unitary$ on the qubits in $\register$ and the
      identity on the qubits in $\qcomp \backslash \register$. Specifically:
      \[
        \unitary_{\register} \coloneqq \SORTgate^\dagger([\register], [\qcomp
        \backslash \register])\ (\unitary \otimes \idmatrix)\
        \SORTgate([\register], [\qcomp \backslash \register])\, ,
      \]
      where $\SORTgate([\register], [\qcomp \backslash \register])$ is the $N
      \times N$ unitary matrix, composed of a series of $\SWAPgate$ gates, that 
      reorders the qubits in $\qcomp$ such that the qubits of $\register$
      precede the qubits of $\qcomp \backslash \register$;
      $\SORTgate^\dagger$ reverses the operation of $\SORTgate$; and 
      $\idmatrix$ is the $2^{n-m} \times 2^{n-m}$ identity matrix.

      Then, the application of $\unitary$ to the qubits in register $\register$ is
      captured in the method $\applyU{\unitary}{\register}[\qcomp]$, defined as
      \begin{align*}
        \applyU{\unitary}{\register}[\qcomp] &: \ket{\qcomp} \longrightarrow
          \unitary_{\register} \ket{\qcomp}\, .
      \end{align*}
    \end{definition}

    \noindent
    Whenever $\qcomp$ is clear from context, we will drop the subscript and
      simply write $\applyU{\unitary}{\register}$. 

    With this definition, it is clear that Alice \emph{may} only be operating on
      the qubits that are available to her, and yet those very same operations may
      be affecting distant parts of $\qcomp$ through entanglement, which
      is now captured in $\applyU{\unitary}{\register}[\qcomp]$ being a
      well-defined evolution of the global pure state $\ket\qcomp$.

    ubsubsection{Interacting entities.}
    With these pieces in place, all that's missing is to define what we mean
      when we say that ``Alice sends a qubit register to Bob''. We will do this
      by letting Alice and Bob each have access to the registers
      $\register[A]$ and $\register[B]$, respectively, which will themselves be
      non-overlapping subregisters of the total system $\qcomp$.

    \begin{definition}[Transferring a subregister]
      Let $\register[A]$ and $\register[B]$ be non-overlapping qubit registers. A
      subregister $\register[R]$ is said to be \emph{transferred} from
      $\register[A]$ to $\register[B]$, denoted $\register[A]
      \apply{\register[R]} \register[B]$, if the elements of $\register[R]$ are
      simultaneously removed from $\register[A]$ and added to $\register[B]$:
      \[
        \register[A] \apply{\register[R]} \register[B]
        \quad \Longleftrightarrow \quad 
        \register[A] \gets \register[A] \backslash \register[R]\, , \quad \register[B] \listapp \register[R]\, .
      \]
    \end{definition}

    \noindent
    Then, Alice sending register $\register[R]$ to Bob simply amounts to
      transferring $\register[R]$ from $\register[A]$ to $\register[B]$.

    Finally, we have reached our goal: A quantum pseudocode rich enough 
      to draw up well-defined interaction diagrams, like the one shown in
      \autoref{fig:aliceandbob}.

    \begin{figure}[t]
      \centering
      \def\arrowlength{3cm}
      \fbox{
        \pseudocode{
          \underline{\text{Alice$(\register[A])$}}\
                \<\< \underline{\text{Bob$(\register[B], \unitary)$}} \\
          (\qubit_0, \ldots, \qubit_{m-1}) \gets \register[A] \\
          \register[R] \gets (\qubit_{i_1}, \ldots \qubit_{i_\ell}) 
            \<\sendmessageright*[\arrowlength]{\register[R]} \\
            \<\< \applyU{\unitary}{\register[R]} \\
            \< \sendmessageleft*[\arrowlength]{\register[R]} \\
          x \sample \meas{}{\register[R]}
        }
      }
      \caption{
        \label{fig:aliceandbob} % *Always* put fig labels in caption, or else autoref gets confused
        Alice sends register $\register[R]$ to Bob, who applies unitary
        $\unitary$ to it and sends it back to Alice,
        who finally measures it and returns the observed value.
        The registers $\register[A]$ and $\register[B]$
        together make up the global quantum system $\qcomp$, left
        implicit in this figure.
      }
    \end{figure}

  \subsection{Measuring Registers}
    \hhnote{
      Split first part away and add it to the primer again, as
      ``Formalizing Measurements''?
    }
    \label{sec:measurements}
      Well---not quite. There is still one piece of formalism that needs to be
        dealt with before we can call it a day: Observations. Our goal is to make the $\meas{R}$
        method in \autoref{fig:aliceandbob} well defined, in much the same
        manner as we did with $\applyU{\unitary}{\register[R]}$---but first we
        will need to take a step back: For, if observations are not 
        unitary operations, what \emph{are} they?
     
    \subsubsection{Projectors.} 
      The answer may be gleaned from the fact, stated earlier, that \emph{upon
        observing a state, the quantum state collapses to the observed state}.
        This means that if you observe the state twice in a row
        (using the same measurement basis, to be defined shortly), then the
        outcome will be the same.
        The state has already collapsed; it can't collapse \emph{more}.

      If we call such an operation $\projector$, then we can write this insight as
        \[ 
          \projector \ket{\psi} = \projector (\projector \ket{\psi}) =
          \projector^2 \ket{\psi} 
          \quad \Longleftrightarrow \quad 
          \projector^2 = \projector\, .
        \]

      \noindent
      Such operators are known as \emph{projectors}, so named because they have
        the effect of \emph{projecting} a higher dimensional object onto a
        lower-dimensional hyperplane---or, in the case of global observations
        (i.e.\ measuring every qubit at once),
        all the way down to one of the basis states.

      For single-qubit states, it is easy to verify that the following
        operators satisfy this definition:

      \[
        \projector_0 \coloneqq \ket{0}\bra{0}
        = \begin{bmatrix} 1 \\ 0 \end{bmatrix} \begin{bmatrix} 1 & 0 \end{bmatrix}
        = \begin{bmatrix} 1 & 0 \\ 0 & 0 \end{bmatrix}\, , \quad
        \projector_1 \coloneqq \ket{1}\bra{1}
        = \begin{bmatrix} 0 \\ 1 \end{bmatrix} \begin{bmatrix} 0 & 1 \end{bmatrix}
        = \begin{bmatrix} 0 & 0 \\ 0 & 1 \end{bmatrix}\, . \quad
      \]

      \noindent
      From the matrix representations above, we see that effect of applying
        $\projector_0$ to a single-qubit state is to \emph{single out} the
        $\ket 0$ state from the superposition: 

      \[
        \projector_0 \ket{\psi} = \projector_0(\alpha \ket{0} + \beta \ket{1})
        = \alpha \ket{0}\, ,
      \]

      \noindent
      and applying $\projector_0$ on it again thereafter
        will have no effect, while applying $\projector_1$ will simply yield
        the zero vector. What's more, we see that 
        $\projector_0 + \projector_1 = \idmatrix[2]$. 
        Clearly, these projectors will be \emph{orthogonal} in the sense that they
        satisfy $\projector_i \projector_j = \delta_{ij} \projector_i$.

      On the other hand, these matrices are manifestly \emph{not} unitary, as
        seen from the fact that they are not even invertible. This reflects the
        irreversible loss of information induced by observation: After
        measuring a state, there is no way to recover the full pre-collapse
        state, other than to re-create it from scratch!

      Let us next define the projector $\projector_\Sspace$ onto some
        subspace $\Sspace$ of the full $N$-dimensional Hilbert space, as
        specified by a subset of the computational basis states
        $\{\ket{x_i}\}_{\forall i \in [k]}$ (where $k \leq N$). This will
        simply be the sum of the relevant basis projectors:

      \[
        \projector_\Sspace = \sum_{\forall i \in [k]} \ket{x_i} \bra{x_i} = 
        \sum_{\forall i \in [k]} \projector_{x_i}\, .
      \]

      \noindent
      The larger the sum, the ``gentler'' the measurement.

      \begin{example}\label{ex:projector}
        Let $\ket\psi = \sum_{i \in [8]} \alpha_i \ket{i}$ be a
        three-qubit state. Then the projector $\projector_{00b}$, given
        by 
        \[
          \projector_{00b} = \ket{000}\bra{000} + \ket{001}\bra{001}
        \] 
        ``selects'' the states 
        in the superposition in which the first and second qubits are zero:
        \[
          \projector_{00b} \ket\psi = \alpha_0 \ket{000} + \alpha_1 \ket{001}\, .
        \]
        We see that $\ket\psi$ has been \emph{projected} from 
        $8$-dimensional Hilbert space down to the $2$-dimensional subspace
        spanned by $\ket{000}$ and $\ket{001}$.
      \end{example}

      \noindent
      Unlike the basis-state projectors $\projector_i$, two general projectors
        $\projector_{\Sspace}$, $\projector_\Tspace$
        need \emph{not} be orthogonal, as subspaces $\Sspace$ and $\Tspace$ may
        be spanned by an overlapping set of basis vectors.

      \begin{exercise}
        Let $\projector_{0b0}$ be given by
        \[
          \projector_{0b0} = \ket{000}\bra{000} + \ket{010}\bra{010}\, ,
        \] 
        and show that $\projector_{0b0}$ and $\projector_{00b}$ are non-orthogonal,
        where $\projector_{00b}$ is as given in \autoref{ex:projector}.
      \end{exercise}

      \noindent
      We will for the most part be interested in sets of orthogonal projectors.
        From the perspective of the computational basis states $\ket x$, this just means
        that no two projectors ``accept'' the same string $x$ as part of their
        respective selections.
        %substrings of $x$ are a particular value, and that no two projectors
        %in the set checks the same bit for the same value. 
        For measurements,
        this translates to no measurement outcome being classified as ``undecided''.

      Projectors may not be unitary, but it is easy to see that they are
        \emph{Hermitian}, meaning
        that they satisfy $\projector^\dagger = \projector$. For the
        computational basis projectors, this follows trivially from the
        matrices being real and symmetric, but it also holds for general
        choices of basis states: Let $\projector_\varphi$ be the projector that
        projects a state $\ket \psi$ onto the one-dimensional subspace spanned
        by the state $\ket\varphi$. Then,
      \[
        \projector^\dagger_\varphi = (\ket{\varphi} \bra{\varphi})^\dagger 
        = \ket{\varphi} \bra{\varphi} = \projector_\varphi\, .
      \]

      \noindent
      Finally, a set of projectors $\{\projector_{\Sspace_i}\}_{\forall i \in [k]}$ is said to be
        \emph{complete} if $\sum_{\forall i \in [k]} \projector_{\Sspace_i} = \idmatrix$.
        In particular, the set of computational-basis projectors $\{
        \projector_i \}_{i \in [N]} = \{ \ket{i}\bra{i} \}_{i \in [N]}$
        is easily seen to be complete: Each element in the sum will simply add one $1$ to
        a previously-zero spot on the diagonal, and zero everywhere else; and
        the sum has $N$ elements, so the full diagonal is covered.
        In the context of measurements, completeness reflects that outcome
        probabilities should sum to one.

      \begin{exercise}
        Provide a set of projectors such that, when taken
        together with $\projector_{00b}$ from \autoref{ex:projector}, the set
        is orthogonal and complete.
      \end{exercise}

    \subsubsection{Projective measurements.}
      However, $\projector \ket{\psi}$ will not itself be a quantum
        state, since it will not have unit norm (unless $\projector =
        \idmatrix$, that is). Additionally, we have not yet
        touched on the concept of measurement \emph{outcomes}. What are the
        values that are actually observed? Something more is needed.
        %This leads us to the concept of a \emph{measurement basis}.

      Let's say that the outcome of some measurement is $m_x$, which we
        interpret as the string $x$. We can formalize this by letting
        $\measurement$ be a \emph{measurement operator} such that
        \[
          \ket\psi \apply{M} m_x \ket x\, ,
        \]
        for one of several possible strings $x$, via some as-yet-undefined
        process. We can also make it more general, by checking which of several
        subspaces $\ket\psi$ belongs to
        \[
          \ket\psi \apply{M} m_{\Sspace_i} \ket{\psi_{\Sspace_i}}\, ,
        \]
        where $\ket{\psi_{\Sspace_i}}$ is the state that results from projecting
        $\ket\psi$ onto subspace $\Sspace_i$.

        %i.e., such that the state collapses to the basis state $\ket\psi$.
      Note that this measurement operator is \emph{not} a projector, since
        repeated applications of $\measurement$ will give back increasing
        powers of $m_x$. However, provided that the measurement
        outcomes are all \emph{orthogonal} (again meaning that none of the
        possible outcomes would be classified as ``undecided''), we are
        guaranteed that there exists a complete set of orthogonal
        projectors $\{\projector_{\Sspace_i}\}_{\forall i \in [k]}$ such that 
        \[
          \measurement = \sum_{\forall i \in [k]} m_{\Sspace_i} \projector_{\Sspace_i}\, .
        \]
        In linear algebra terms, this defines the \emph{spectral decomposition} of the
          operator $\measurement$, with the eigenvalues $m_{\Sspace_i}$.

        In the case that you are just reading off the values of $n$ qubits
          directly in what we'll refer to as a ``computational-basis
          measurement'', this will be
          \[
            \measurement = \sum_{\forall i \in [N]} m_i \projector_i
            = \sum_{\forall i \in [N]} m_i \ket{i} \bra{i}\, ,
          \]
          which, when expressed in the computational basis, is just a sum of
            $N$ matrices with the single non-zero element $m_i$ in the $i$'th
            position along its diagonal.

        Finally, then, we are ready to formally define our measurement procedure.

        \begin{definition}[Projective measurement]
          Let $i$ be an $n$-length bitstring, let $\{ \projector_{\Sspace_i}
            \}_{\forall i \in [k]}$ be a complete, orthogonal set of projectors
            (projecting onto the orthogonal subspaces $\Sspace_i$ of the full Hilbert space),
            and let $\measurement \coloneqq \sum_{\forall i \in [k]} m_{\Sspace_i} \projector_{\Sspace_i}$
            be a projective measurement operator.
            %
            Then, \emph{measuring} a state $\ket{\psi}$ under $\measurement$ yields outcome
            $\Sspace_i$ with probability:
            \[
              \prob{\abs{\measurement \ket\psi} = m_{\Sspace_i}}
              = \abs{\projector_{\Sspace_i} \ket{\psi}}^2\, ,
            \]
            and, conditioned on $\Sspace_i$ being observed, the act of measuring 
            transforms the state as follows:
            \[
              \ket{\psi} \apply{\measurement} \frac{\projector_{\Sspace_i}
              \ket{\psi}}{\abs{\projector_{\Sspace_i} \ket{\psi}}}\, .
            \]
            The value $\projector_{\Sspace_i} \ket{\psi} \coloneqq \alpha_{\Sspace_i} \in \CC$
            is known as the \emph{amplitude} of observing outcome $\Sspace_i$.
        \end{definition}
        
        \noindent
        Again specializing to the case of measuring $n$ qubits in the
          computational basis and reading off the resulting string, we see that
          the probability of observing any particular string $x$ is given by
          \[
            \prob{\abs{\measurement \ket\psi} = m_x} = \abs{\projector_x \ket{\psi}}^2
            = \abs{\alpha_x \ket x}^2 = \abs{\alpha_x}^2\, ,
          \]
          with the post-measurement state 
          \[
            \ket\psi \apply{\measurement} 
            \frac{\projector_x \ket{\psi}}{\abs{\projector_x \ket{\psi}}}
            = \frac{\alpha_x \ket x}{\abs{\alpha_x \ket x}}
            = \frac{\alpha_x }{\abs{\alpha_x}} \ket x
            \longrightarrow \ket x\, ,
          \]
          as required (where the final step does away with the unobservable
          global phase $\frac{\alpha_x }{\abs{\alpha_x}} = e^{i\theta}$).
        
        You may wonder, what \emph{are} the values of the measurement outcomes, $m_x$?  Let
          us look at single-bit measurements: $m_0$ will signal outcome $\ket 0$, and
          $m_1$ will signal outcome $\ket 1$. Though tempting, we surely can't set
          $m_b = b$, since the outcome $\ket 0$ would then multiply away the
          entire state! Instead, the canonical choice\footnote{
            Historically, this choice corresponds to the possible outcomes of
            the \emph{spin} observable for fermionic particles, such as
            electrons, with the value $+1$ corresponding to
            ``spin-up'' (the spin vector points in the positive direction of
            the $z$-coordinate) and the value $-1$ meaning ``spin-down'' (the
            spin vector points in the negative $z$-direction).
          } is to let $m_0 = 1$ and $m_1 = -1$. Then, 

        \[
          \measurement = m_0 \projector_0 + m_1 \projector_1
          = \begin{bmatrix} 1 & 0 \\ 0 & -1 \end{bmatrix}\, .
        \]
        You may recognize this matrix from earlier: It's the $\Zgate$ gate! In
          fact, computational-basis measurements are often referred to as
          ``$\Zgate$ measurements'' for this reason. As you can see, this
          choice also has the useful property that the measurement operator is
          unitary, and so may be identified with an operator, with the possible
          outcomes represented by the eigenvalues and eigenstates of that operator.

        This mapping, from a unitary to a measurement, is possible to do for any
          operator that is both unitary and Hermitian, meaning that 
          $\unitary^{-1} = \unitary^\dagger = \unitary$. Such unitaries are
          easy to recognize: They are the ones
          with real values along the diagonal, and complex-conjugate symmetry
          over the off-diagonals: $\unitary{[i,j]} = \unitary{[j,i]}^*$. 

        For example, in addition to the $\Zgate$ and $\Xgate$, we also
          have a $\Ygate \coloneqq i\Xgate\Zgate$. It turns out that any
          measurement that is possible to perform on a single qubit, can be
          written as a linear combination of these three matrices, which are
          known as the \emph{Pauli matrices} after their discoverer.

        \begin{exercise}
          Show that interpreting the $\Xgate$ (a.k.a. ``NOT'') gate as a
          measurement operator leads to
          the possible outcomes $\ket{+}$ and $\ket{-}$ (where these are as
          defined in \autoref{ex:HadamardBasis}). 
          What are the values of $m_+$ and $m_-$?
        \end{exercise}

        \begin{exercise}
          Write out the matrix representation of $\Ygate$ in the computational
          basis, and find the possible measurement outcomes it gives rise to
          when interpreted as a measurement operator by identifying its
          eigenstates and eigenvalues.
        \end{exercise}

      \subsubsection{Measuring registers.}
        We now give the last piece of the puzzle, namely a pseudocode for
          perfoming measurements on registers that may or may not be entangled
          with other parts of the system (i.e.\ may or may not be \emph{pure}).

        \begin{definition}[Operating on registers]\label{def:apply}
          Let $\qcomp$ be an $n$-qubit quantum memory in quantum state
          $\ket{\qcomp}$, let the register $\register$ be an $m$-qubit
          subsystem of $\qcomp$, and let $\unitary$ be an
          $M \times M$ unitary matrix. 

          Next, define $\unitary_{\register}$ to be the $N \times N$ unitary matrix that
          has the effect of applying $\unitary$ on the qubits in $\register$ and the
          identity on the qubits in $\qcomp \backslash \register$. Specifically:
          \[
            \unitary_{\register} \coloneqq \SORTgate^\dagger([\register], [\qcomp
            \backslash \register])\ (\unitary \otimes \idmatrix)\
            \SORTgate([\register], [\qcomp \backslash \register])\, ,
          \]
          where $\SORTgate([\register], [\qcomp \backslash \register])$ is the $N
          \times N$ unitary matrix, composed of a series of $\SWAPgate$ gates, that 
          reorders the qubits in $\qcomp$ such that the qubits of $\register$
          precede the qubits of $\qcomp \backslash \register$;
          $\SORTgate^\dagger$ reverses the operation of $\SORTgate$; and 
          $\idmatrix$ is the $2^{n-m} \times 2^{n-m}$ identity matrix.

          Then, the application of $\unitary$ to the qubits in register $\register$ is
          captured in the method $\applyU{\unitary}{\register}[\qcomp]$, defined as
          \begin{align*}
            \applyU{\unitary}{\register}[\qcomp] &: \ket{\qcomp} \longrightarrow
              \unitary_{\register} \ket{\qcomp}\, .
          \end{align*}
        \end{definition}

        \begin{definition}[Measuring a register]\label{def:measure}
          Let $\qcomp$ be an $n$-qubit quantum memory in pure state
            $\ket{\qcomp}$, let the register $\register$ be an $m$-qubit
            subsystem of $\qcomp$, and let $\measurement$ be an $m$-qubit
            projective measurement that check the state for membership in
            Hilbert subspaces $\{ \Sspace_i \}$.

          Next, define $\measurement_{\register}$ to be the $n$-qubit
            projective measurement that
            has the effect of applying $\measurement$ on the qubits in $\register$ and the
            identity on the qubits in $\qcomp \backslash \register$
            (constructed in the same manner as $\unitary_{\register}$ in
            \autoref{def:apply}).

          Then, performing measurement $\measurement$ on the qubits in register
            $\register$ is captured in the (probabilistic) method
            $\meas{\measurement}{\register}[\qcomp]$ which observes
            $\ket{\qcomp}$ under $\measurement_{\register}$. Observing outcome
            $i$, designated by the measurement procedure returning value $m_i$,
            is written
          \[
            i \sample \meas{\measurement}{\register}[\qcomp]\, ,
          \]
            and the post-measurement will be 
          \[
            \ket{\qcomp} \apply{\measurement} \ket{\qcomp}_{\Sspace_i}\, ,
          \]
          where $\ket{\qcomp}_{\Sspace_i}$ is the result of projecting the
            qubits of $\register$ onto subspace $\Sspace_i$, and applying the
            identity to the rest.

          We will omit the subscript whenever $\qcomp$ is clear from context,
            and the measurement operator $\measurement$ whenever we are
            performing a computational-basis measurement; for example, measuring
            $\ell$ qubits and observing the bit string $x \in \bin^\ell$ will
            be written simply
          \[
            x \sample \meas{}{\register}\, .
          \]
        \end{definition}

        \noindent
        To showcase the utility of our notation, we can now restate the Born
          rule in a very compact form.

        \begin{definition}[The Born rule]\label{def:Born}
        If qubit register $\register$ is in pure state $\ket{\register} = \sum_i \alpha_i \ket{x_i}$, then
          \[
            \probsample{x \sample \meas{}{\register}}{x = x_i} = \abs{\alpha_i}^2\, .
          \]
        \end{definition}
